\documentclass[a4j,fleqn,10pt]{jarticle}
\usepackage{ipsj}
\usepackage{txfonts}
\usepackage[symbol*,perpage]{footmisc}
\usepackage[dvipdfm]{graphicx}
%\usepackage{slashbox}
%\usepackage{amsmath}
\newcommand{\argmax}{\mathop{\rm arg~max}\limits}
\newcommand{\argmin}{\mathop{\rm arg~min}\limits}

%-------------------------------------------jtitle and jauthors information
\begin{document}

\jtitle{類似貢献度に基づく文書検索結果の可視化分析法と結果の評価法}

\jcontact{
\dag 神奈川大学大学院 理学研究科
}

\jauthor{尾崎　了祐\dag}
%登録番号　5401
\maketitle

%-------------------------------------------etitle and eauthors information
\footnotetext[0]{\hspace*{-5.2mm}Visualization and Evaluation of Document Search Results Based on Similarity Contribution} %英タイトル
\footnotetext[0]{\hspace*{-6.5mm}\dag Ryosuke OZAKI \quad \dag Kazumi SAITO \quad} %著者情報
\footnotetext[0]{\hspace*{-6.5mm}\dag Kanagawa University} %自分の英語版所属

%-------------------------------------------document

\section{はじめに}
ビッグデータ技術の進展により, 膨大な文書データから, 有用な情報を利活用するニーズは益々高まっている~\cite{ikeda17}.
本研究では, 与えられた文書クエリの類似検索結果文書集合を可視化する課題に焦点を当てる. 
ところが, 著者らの知る限り, この課題に関する先行研究は少数で~\cite{oda04}, あまり研究が進んでいない状況である. 
本論文では, 検索結果の各文書をノードと見なし, それぞれを類似貢献度ベクトルに変換する. 
そして, k-最近傍グラフを構築した可視化において, 検索結果文書がクエリに対し, 
どのような説明語で類似するかを付与するとともに, グループ分けして配色する新手法を提案する. 
livedoor ニュースコーパスを用いた評価実験では, 可視化例に基づく定性評価とともに, 
説明語となる単語数の定量評価により, 提案法の有効性を検証する.
さらに, 本論文では, 各種可視化法の有効性を, エッジ長の変動係数を用いて定量評価する方法を提案する. 
livedoorニュースコーパスを用いた評価実験では, ４つの手法でグラフを可視化し, 
定量評価と可視化例に基づく定性評価を比較することで, 定量評価法の妥当性を検証する.

\section{先行研究}
大規模文書検索結果のクラスタリングと可視化を組み合わせている先行研究として,~\cite{oda04}が挙げられる.
当研究は, 検索結果文書全体を対象として, 文書間の類似度を考慮した可視化結果を生成している点で共通しているが,
文書のクラスタリング手法として, 文書間の距離ベクトルを最長距離法で算出する階層型クラスタリングを採用しており,
検索結果文書を3次元空間上で可視化している.
一方で, 本研究では文書間の類似ベクトルをコサイン類似度で算出し,
そのベクトルをクエリを考慮した類似貢献度ベクトルに変換する,非階層型クラスタリングを採用している. 
検索結果文書も, 2次元上で可視化している点で異なっている.
グラフ可視化手法を対象とした先行研究~\cite{saito},~\cite{YST},~\cite{BETT},~\cite{C}では, 
２次元の可視化結果を生成し, ネットワーク分析の有効性を実証している.

\section{類似貢献度を用いた提案手法}
総数 $N$ の文書集合を ${\cal N} = \{1, \cdots, n, \cdots, N \}$ とし, 
語彙数 $D$ の出現単語集合を ${\cal D} = \{1, \cdots, d, \cdots, D \}$ とする. 
また, 各文書 $n$ とクエリ文書の $D$-次元特徴ベクトルを ${\bf x}_n$ と ${\bf q}$ とし,  
文書間でのコサイン類似度を求めるため, 特徴ベクトルはノルムが $1$ に正規化されているとする. 
  
提案可視化分析法の手順を以下に述べる. 
まず, クエリ ${\bf q}$ に対し, 文書集合 ${\cal N}$ から $h$-最類似文書集合 ${\cal R}(h) \subset {\cal N}$ を求める. 
次に, 類似文書 $n \in {\cal R}(h)$ に対し, 特徴ベクトル ${\bf x}_n$ から単語 $d$ の類似貢献度を
\begin{eqnarray}
y_{n,d} & = & \frac{q_d \times x_{n,d}}{\sum_{d' \in {\cal D}} q_{d'} \times x_{n,d'}} 
\label{sc1}
\end{eqnarray}
で求め, $D$-次元類似貢献度ベクトル ${\bf y}_n$ を構築する. 
ここで, $y_{n,1} + \cdots + y_{n,D} = 1$ となる. 
そして, 類似貢献度ベクトル ${\bf y}_n$ 間の距離に基づき, ${\cal R}(h)$ をノード集合として 
$k$-最近傍 (NN: Nearest Neighnour) グラフ $G({\cal R}(h), {\cal E})$ を構築する. 
ここで, ${\cal E}$ はノード間に張られたリンク集合を表す. 
最後に, 単語 $d$ で類似貢献度が最大となるノード(文書)集合を
\begin{eqnarray}
{\cal C}(d) & = & \{ n \in {\cal R}(h)~|~d = \argmax_{d' \in {\cal D}} (y_{n,d'}) \} 
\label{sc2}
\end{eqnarray}
として求め, そこに属すノード群を同色にし, 単語 $d$ を説明語として表示可能にした 
$k$-最近傍 $G({\cal R}(h), {\cal E})$ の可視化結果を出力する. 

比較評価のため, 式~(\ref{sc1})でクエリとの類似貢献度を考慮せず, 
類似文書 $n \in {\cal R}(h)$ それぞれの類似度として, $z_{n,d} = x_{n,d}^2$を求め, 
$D$-次元ベクトル ${\bf z}_n$ を構築し, 説明語を $\argmax (z_{n,d'})で$求める方法を比較法と呼ぶ. 
なお, 先行研究~\cite{oda04}では, クエリを考慮した類似貢献度ベクトルを用いず, 
検索結果文書間の類似度を求める点で比較法の範疇と見なせる. 
また, 提案法において, グラフ構築に最小全域木(MST: Minimum Spanning Tree)を用いる方法をMST法と呼び比較評価する.

\section{類似貢献度の効果検証}
本実験では，文書数 $N = 7,367$ の livedoor ニュースコーパスを文書集合 ${\cal N}$ とし, 
MeCab~\cite{kudo04}で出力される形態素をすべて単語として扱い, 語彙数は $D = 89,993$ とした. 
なお, グラフ可視化には, ばねモデル~\cite{kamada89}を採用した.
まず図~\ref{fig0}に，提案法, 比較法, 及び, MST法での可視化結果の例を示す. 
ここで比較的大きな赤丸で描いたクエリは, 文書集合 ${\cal N}$ から随意に選択した, 
ゆるキャラ「カクカクシカジカ」によるダイハツ車CMに関する文書である. 
そして, 類似文書数を $h = 100$ に設定し ${\cal R}(h)$ を求めた.  
ノード配色については, 類似文書 ${\cal R}(h)$ を類似度で降順に並べ, 
同じ説明語となるノード群は同色となるよう, CMYKカラーシステム24色をサイクリックに割り当てた. 

図~\ref{fig0}(a)に示す提案法での可視化結果からは,  
CM, 店, クルマを説明語とし, 比較的密結合する多数のノード群とともに,  
比較的大きな赤丸のクエリ近くに, 少数ながらダイハツ, ゆるキャラを説明語とするノード群も見て取れる. 
その他の単一説明語としては, 自動車, キャンペーンなどが出現する. 
これら結果より, ある観点で類似する文書数がどの程度かなど視覚的に把握可能と言える. 
これに対し, 図~\ref{fig0}(b)に示す比較法での可視化結果では, 
殆どのノード結合は異なる説明語 (配色) となり, 全体的な類似構造把握の点などでも限界が見受けられる. 
一方, 図~\ref{fig0}(c)に示すMST法での可視化結果と比較すれば, 
ノード配色は提案法と同一であるが, 密結合する部分構造把握の点などで, 提案法に望ましい性質が見て取れる.

\begin{figure}[t]
   \begin{center}
    \includegraphics[width=\hsize]{figs/fig2a.png}
    \mbox{(a)提案法}
   \end{center}
   \begin{center}
    \includegraphics[width=\hsize]{figs/fig2b.png} 
    \mbox{(b)比較法}
   \end{center}
   \begin{center}
    \includegraphics[width=\hsize]{figs/fig2c.png}
    \mbox{(c) MST法}
   \end{center}
 \caption{可視化結果の例}
 \label{fig0}
\end{figure}

次に, 類似貢献度ベクトルを用いる効果を定量評価する. 
まず, $h$-最類似文書 ${\cal R}(h)$ における各単語 $d$ での類似貢献度の期待値を $u_d = \sum_{n \in {\cal R}(h)} y_{n,d}/h$ で求め, 
これら値を降順 $u_{d(r)} \geq u_{d(r+1)}$ に並び替え, その累積値を $v_s = \sum_{r=1}^s u_{d(r)}$ で求める. よって, $v_D = 1$ となる. 
図~\ref{fig1}に，クエリ文書を文書集合 ${\cal D}$ のそれぞれに設定したときの累積値 $v_s$ の期待値を示す. 
ここで, p050, p100, p200 は上述した提案法で検索文書数を $h = 50, 100, 200$ に設定した結果であり, 
c050, c100, c200 には, 類似貢献度期待値を特徴量での期待値 $u_d' = \sum_{n \in {\cal R}(h)} z_{n,d}/h$ 
に置き換えた比較法において $h = 50, 100, 200$ に設定した結果である. 
図~\ref{fig1}より，検索文書数 $h$ に依らず, 提案法では上位 $r = 10$ 単語程度で累積貢献値が0.5程度になるのに対して, 
比較法では, そのために上位数百単語程度が必要になる. 
すなわち, 提案法では, 比較的少数の単語に類似貢献度が集中するので, 妥当なノード配色や説明語選択の実現が期待できる. 

\begin{figure}[t]
  \begin{center}
  \includegraphics[width=\hsize]{figs/fig1.png}
  \end{center}
  \caption{貢献度期待値の累積分布}
  \label{fig1}
\end{figure}

\section{可視化グラフの定量的評価法}
任意のノードをホバーすることで, 文書詳細を閲覧できる機能をもつ可視化グラフを想定し, 
与えられた文書クエリの類似検索結果文書集合を, グラフを用いて可視化する手法を提案している~\cite{ozaki1}.
一方で, グラフ可視化には他の手法~\cite{saito},~\cite{kamada89}を用いることも可能だが, 可視化結果を定量評価する方法はあまり知られていない.
本研究では, 各種可視化法の有効性を, エッジ長の変動係数を用いて定量評価する方法を提案している~\cite{ozaki2}.
まず, 提案法を述べる前に, 良い可視化グラフといえる基準を定義する.
本研究では下記項目を設定した.
\begin{enumerate}
    \item 隣接ノードが離れすぎずに描画される
    \item 重なっているノードが存在しない
\end{enumerate}
項目１は, 隣接ノード同士が一定の範囲内で描画されるか評価している.
非常に離れた位置に点在している場合, ホバーしにくいためである.
項目２は, 複数のノードが重なって描画されないか評価している.
重なるとホバーが困難なため, 最悪の場合, 文書詳細を閲覧できないノードが出現する.
従って, 上記項目をどの程度満たすか定量的に評価可能な手法を提案することが目的となる.
次に, 提案可視化評価法の手順を以下に述べる.
エッジ長の集合を ${\cal E} = \{1, \cdots, e, \cdots, E \}$ とし,
${\cal E}$ に対し, 平均 $\mu$, 標準偏差 $\sigma$ を求める.
次に, 変動係数を
\begin{equation}
  {cv} = \frac{\sigma}{\mu} \times 100
  \label{sc3}
\end{equation}
で導出する.
式(\ref{sc3})は, エッジ長のばらつき度合いを示す数値であり,
この値が小さい程, 各エッジの長さは全体的に同程度となる.
故に, 隣接ノードが一定の範囲内に描画される割合が高くなる.
逆に, 値が大きい程, 各エッジの長さに「ばらつき」が生じる.
故に, 短すぎる, 長すぎる等の不自然なエッジ長をもつ割合が大きくなり, 
結果的に隣接ノード同士の距離が不自然なペアが出現する.
従って, 変動係数が小さくなる程, 項目１を高水準で満たすグラフとなるため,
式(\ref{sc3})を用いることで, ある程度の定量評価が可能になる.

次に, 項目２に該当するケースも考慮する.
ここでは, 半径 $r$ の円として可視化する.
例えば, エッジの生成範囲が一方向に偏っている場合, 
ノードが重なる割合が高くなり, ホバーが困難なグラフが生成されてしまう.
従って, ノード集合を ${\cal N} = \{1, \cdots, n, \cdots, N \}$, 
ノード $n$ の座標を $(X_n, Y_n)$, $\forall n \neq m$ の条件で重なっているノードの個数を $on$とし,
\begin{equation}
  {on} = |\left\{ n, m \mid \sqrt{(X_n - X_m)^2 + (Y_n - Y_m)^2} \leq 2r \right\}|
  \label{sc4}
\end{equation}
を求め, 定量評価に用いる.
この値が小さくなる程, 項目２を高水準で満たすグラフとなる.

最後に, 可視化ネットワーク数を$L$, 可視化手法を$M$とした$L*M$個のデータを用い, 
X軸には評価項目１を式(\ref{sc3})で求めた値を,
Y軸には評価項目２を式(\ref{sc4})で求めた値を設定した散布図を生成する. 

\section{定量評価法の効果検証}
提案法での評価項目１,２を用いて, 可視化グラフを定量評価可能か実験する.
可視化手法$M$には, スペクトラル法(se), 多次元尺度法(mds),
クロスエントロピー法(ce), ばねモデル法(kk)を採用した.
また, 定量評価する際の実験データとして, livedoor ニュースコーパスから文書検索を行い, 
ノード数を $100$ に設定した全7367件の可視化グラフを生成した. 
本稿では, 標本集団として抽出した$L=8$, $M=4$の範囲内で図~\ref{fig1}の散布図を用いて説明する.
X軸には評価項目１を式(\ref{sc3})で求めた値を,
Y軸には評価項目２を式(\ref{sc4})で求めた値を設定した.
左下に近い位置にプロットされるサークルである程, 項目１,２を高水準で満たすグラフとなる.
また, 各サークルの数値はクエリ文書の $id$ を示しており, 同番のサークルは同様のクエリ文書を用いて可視化グラフを生成している.
先ほどの可視化例で示した~\ref{fig0}は, クエリ文書$1$である.

図~\ref{fig2}より, ばねモデル法はどれも左下にプロットされており, 可視化グラフとして高水準である.
一方で, スペクトラル法は全体的に右上にプロットされているため低水準と判断できる.
多次元尺度法, クロスエントロピー法は全体的に図の中央に位置しており, 
それぞれを比較してみると, 項目１は多次元尺度法, 項目２はクロスエントロピー法に分があることが見て取れる.

次に, 実際に生成された可視化グラフを確認し, 先ほどの定量評価と結果が一致するか確認する. 
今回は, 各種可視化グラフの評価値が全体的に同程度となった,
すなわち, 目視での評価が困難である, クエリ文書$6$を対象として確認した.

クエリ文書$6$のグラフ可視化結果を図~\ref{fig3}に示す. 
スペクトラル法(a), 多次元尺度法(b), クロスエントロピー法(c)では,
離れている, 部分的に重なっている隣接ノードが存在するが,
ばねモデル法(d)では全く存在しないことが見て取れる. 
これら結果より, 定性評価でも, ばねモデル法が高水準であり, 定量評価と結果が一致した.

\begin{figure}[t]
  \begin{center}
  \includegraphics[width=\hsize]{figs/ipsj2025/fig1.png}
  \end{center}
  \caption{可視化グラフの評価値を示す散布図}
  \label{fig2}
\end{figure}

\begin{figure}[t]
   \begin{center}
    \includegraphics[width=\hsize]{figs/ipsj2025/fig2a.png} 
    \mbox{(a)スペクトラル法}
   \end{center}
   \begin{center}
    \includegraphics[width=\hsize]{figs/ipsj2025/fig2b.png}
    \mbox{(b) 多次元尺度法}
   \end{center}
   \begin{center}
    \includegraphics[width=\hsize]{figs/ipsj2025/fig2c.png}
    \mbox{(c) クロスエントロピー法}
   \end{center}
   \begin{center}
    \includegraphics[width=\hsize]{figs/ipsj2025/fig2d.png}
    \mbox{(d)ばねモデル法}
   \end{center}
 \caption{クエリ文書$6$の各種グラフ可視化結果}
 \label{fig3}
\end{figure}

\section{おわりに}
本論文では, 検索結果の各文書をノードと見なし, 
それぞれを類似貢献度ベクトルに変換し, 
k-最近傍グラフを構築した可視化において, 
検索結果文書がクエリに対し, どのような説明語で類似するかを付与するとともに, 
グループ分けして配色する新手法のグラフ可視化法を提案した. 

また, 可視化グラフの評価方法として, エッジ長の変動係数を定量評価する手法を採用することで, 
どの手法を用いた可視化グラフが望ましいのかを客観的に評価可能となり, 
検証実験では, 定性評価での結果と一致することを確認し, 提案法の妥当性を検証した. 
今後の研究では，提案法の妥当性を様々な可視化手法で検証する．

%\begin{description}
%\item[{\bf 謝辞}]
%\small 本研究は，科学研究費補助金基盤研究(C)(No．26330138)の補助を受けた．
%\end{description}

\fontsize{9.5pt}{0pt}\selectfont
\begin{thebibliography}{10}
\bibitem{ozaki1}
尾崎 了祐, 斉藤和己.
類似貢献度に基づく文書検索結果の可視化分析法.
情報処理学会第８６回全国大会, 2024. 
\bibitem{ozaki2}
尾崎 了祐, 斉藤和己.
類似貢献度に基づく文書検索結果の可視化評価法.
情報処理学会第８７回全国大会, 2025. 
\bibitem{saito}
斉藤和巳.
ネットワークの可視化技術―大規模情報からの意味情報の抽出―.
数理科学, 536, 78-83, 2008.
\bibitem{ikeda17}
池田 弘行.
機械学習を用いた 文章データの解析・可視化技術.
東芝レビュー, 5, 68-69, 2017. 
\bibitem{oda04}
小田 良治, 土井 章男.
大規模文書検索結果のクラスタリングと可視化.
情報処理学会第66回全国大会, 2004. 
\bibitem{kudo04}
T. Kudo, K. Yamamoto, and Y. Matsumoto, 
Applying Conditional Random Fields to Japanese Morphological Analysis, 
Proceedings of the 2004 Conference on Empirical Methods in Natural Language Processing (EMNLP-2004), 230-237, 2004.
\bibitem{kamada89}
T. Kamada and S. Kawai. 
An algorithm fordrawing general undirected graphs. 
Information Processing Letters. 31, 7-15, 1989.
\bibitem{BETT}
G. Battista, P. Eades, R. Tamassia, and I. Tollis, 
Graph drawing: An annotated bibliography Prentice-Hall, New Jersy, 1999.
\bibitem{C}
F. R. K. Chung,
Spectral Graph Theory, Number 92 in CBMS Regional Conference Series 
in Mathematics American Mathematical Society, 1997. 
\bibitem{YST}
山田 武士, 斉藤 和巳, 上田 修功, クロスエントロピー最小化に基づくネットワークデータの埋め込み.
情報処理学会論文誌, 44 2401-2408, 2003.

\end{thebibliography}

\end{document}
